\section{Introduction}

An \(\qc{n,k,d}\) quantum code uses \(n\) physical qubits to store
\(k\) logical qubits. Its distance \(d\) is the smallest number of
qubits on which a Pauli operator can act nontrivially on the encoded
information while commuting with every stabilizer. Distance measures
the weight of the smallest undetected logical error. A code of distance
\(d\) corrects arbitrary errors on up to \(\lfloor(d-1)/2\rfloor\)
qubits. Both distance five and distance six correct two-qubit errors;
distance six also detects errors of weight five.

For fixed numbers of physical and logical qubits, determining the largest
possible minimum distance is a basic problem in quantum coding theory.
Constructions give lower bounds, while weight-enumerator identities and
other general bounds give upper bounds. When these differ by one, the
remaining existence question may depend on structure that an ordinary
weight enumerator does not record.

Let \(d_{\max}(n,k)\) denote the largest minimum distance of a binary
stabilizer code encoding \(k\) logical qubits into \(n\) physical qubits.
We prove the following nonexistence result.
\begin{theorem}\label{thm:main}
There is no binary stabilizer code with parameters
\(\qc{17,3,d\geq6}\). Hence
\[
 d_{\max}(17,3)=5.
\]
The conclusion holds for both pure and degenerate stabilizer codes.
\end{theorem}
The matching \(\qc{17,3,5}\) construction is listed in Grassl's tables
\cite{grassl}; it extends a \(\qc{15,3,5}\) code by two coordinates.
Its distance is checked by complete centralizer enumeration
(Section~\ref{sec:verification}). The exclusion closes the distance
interval \(5\)--\(6\) at \((n,k)=(17,3)\).
The same method excludes \(\qc{16,5,5}\); that proof is given separately.

\subsection{Outline of the method}

The proof begins by forcing a stabilizer word whose support leaves few
enough complementary coordinates for a complete classification. Let
\(T\) be that support and \(R\) the entire stabilizer supported on
\(T\). To shorten the centralizer on \(T\), retain only its words that
are identity on every coordinate of \(T\), then delete those coordinates.
Call the resulting additive code \(C\).
Every centralizer word has a component on \(T\) and a component on its
complement. Once the first component is fixed modulo \(R\), the second
belongs to a single coset of \(C\). The possible codes can be
described by linear maps between quotient spaces.

\begin{figure}[htbp]
\centering
\begin{tikzpicture}[>=Stealth,
 proofbox/.style={draw=black!60,fill=black!4,rounded corners=1pt,
 text width=4.25cm,minimum height=1.05cm,align=center,font=\small},
 node distance=0.65cm and 0.7cm]
\node[proofbox] (assume) {Assume the code exists};
\node[proofbox,right=of assume] (force) {Force a weight-eight stabilizer word};
\node[proofbox,right=of force] (split) {Split the qubits on its support};
\node[proofbox,below=of split] (classify) {Classify the shortened codes};
\node[proofbox,left=of classify] (test) {Test every extension map};
\node[proofbox,left=of test] (none) {No map, so no code};
\draw[->] (assume)--(force);
\draw[->] (force)--(split);
\draw[->] (split)--(classify);
\draw[->] (classify)--(test);
\draw[->] (test)--(none);
\end{tikzpicture}
\caption{The structural reduction and the finite extension test.}
\label{fig:proof-flow}
\end{figure}

A coset of a subspace consists of all vectors obtained by adding its
members to one fixed vector. The quotient is the vector space of these
cosets: it treats two vectors as equivalent when their difference lies
in that subspace. Here this identifies local Pauli words that differ by
a stabilizer supported on the selected coordinates.

Put \(E=R_T^\perp/R\). Write \(c_R(e)\) for the minimum weight of a
representative of \(e\in E\), and write \(\delta(u)\) for the minimum
weight of a coset \(u\) of \(C\). In suitable coordinates, an extension
is described by a binary matrix \(M\) satisfying
\begin{equation}\label{eq:intro-criterion}
 \Gamma=MJM^\tr,\qquad c_R(e)+\delta(Me)\geq d\quad(e\ne0).
\end{equation}
Here \(J\) is the symplectic matrix on \(E\), and \(\Gamma\) is the
Gram matrix of a basis of \(C^\perp\). The first equation is the
commutation condition. The second states that the least weight in each
corresponding pair of cosets is at least \(d\).
Theorem~\ref{thm:extension} proves that these conditions are necessary
and sufficient when \(d(C)\geq d\) and every stabilizer word of weight
below \(d\) lies in \(R\); this holds for pure codes.

For example, suppose \(R=\langle ZZZZ\rangle\) on four coordinates.
Here \(I\) denotes the identity and \(Z\) a nonidentity Pauli letter.
The local class \(ZIII+R=\{ZIII,IZZZ\}\) has minimum weight one.
If its image under \(M\) had coset minimum weight three, the two
parts could be combined into a centralizer word of weight four outside
\(R\). In the setting of Theorem~\ref{thm:extension}, this rules out
distance five. The inequality in \eqref{eq:intro-criterion} makes this
test for every local class at once.

Every nonzero element of \(\ker M\) represents a logical operator
supported entirely on \(T\). The weights and symplectic products of
these operators restrict the hull \(C\cap C^\perp\) of the shortened
code. These restrictions divide the problem into cases with prescribed
kernels, before any extension matrix is enumerated.

For seventeen qubits, exact weight-enumerator certificates show that a
hypothetical code is pure and has a stabilizer word of weight eight.
A quadratic-form argument bounds the binary dimension of the shortened
nine-coordinate code by six. It follows that there are just five
possible local stabilizers \(R\): one of rank one and four of rank two.
The rank-one type has five admissible kernel cases, determined by the
kernel dimension, restricted symplectic rank, and shortened-code hull.
These restrictions reduce the remaining problem to complete
classifications on nine coordinates and finite extension searches.

\subsection{Map of the proof}\label{sec:map}

Table~\ref{tab:map} lists every case of the proof, the argument that
excludes it, the size of the corresponding finite computation, and the
place where it is handled. We use \emph{type} for the local stabilizer
\(R\) up to physical equivalence, labeled by its nonzero weights;
\emph{kernel case} for the invariants \((q,t,h)\) of
Section~\ref{sec:rank-one}; and \emph{weight profile} for the pair
\((A_6,A_7)\) of stabilizer weight counts.

\begin{table}[htbp]
\centering\small
\renewcommand{\arraystretch}{1.12}
\begin{tabularx}{\linewidth}{@{}>{\raggedright\arraybackslash}p{3.15cm}>{\raggedright\arraybackslash}X>{\raggedright\arraybackslash}p{3.55cm}l@{}}\toprule
Case&How it is excluded&Classes or cases&Where\\\midrule
Stabilizer words of weight \(1\)--\(5\)&Split-coset integer certificates; purity follows&5 certificates&Prop.~\ref{prop:17-enumerator}\\
No weight-eight word&Integer certificate with \(A_8=0\)&1 certificate&Prop.~\ref{prop:17-enumerator}\\
Local rank \(r\geq3\)&Nine-coordinate dimension bound&None&Lemma~\ref{lem:nine-bound}\\
Shortened codes&Complete classifications&32,865 classes (dim.~5); 56 classes (dim.~6)&Section~\ref{sec:classification}\\\midrule
Rank two, \((6,6,8)\)&Counting leaves 48 of 56; no \(Z\)-only tuple passes&641 first-image cases&Prop.~\ref{prop:17-mixed}\\
Rank two, \((6,8,8)\)&Counting leaves 52 of 56; searches for each weight profile&1,950 first-image cases&Prop.~\ref{prop:17-mixed}\\
Rank two, \((7,7,8)\)&Counting leaves 52 of 56; extension searches&650 first-image cases&Prop.~\ref{prop:17-mixed}\\
Rank two, \((8,8,8)\), hull \(2,6\)&Hull identities; fixed-kernel searches, \(q=1,3\)&41 classes; 1,375 fixed-kernel problems&Prop.~\ref{prop:17-plane}\\
Rank two, \((8,8,8)\), hull \(0\)&Profile exclusions and extension searches&15 classes; 72,827 first-image cases&Prop.~\ref{prop:17-plane}\\\midrule
Rank one, \(q=4\)&Parity and isotropy&None&Lemma~\ref{lem:qthree}\\
Rank one, \((1,0,1)\), \((2,2,1)\), \((2,0,3)\)&Fixed-kernel extension searches&32,747 classes; 2,419,796 records&Prop.~\ref{prop:17-rankone}\\
Rank one, \((3,2,3)\)&Recentering screen; fixed-kernel searches&5,995 classes; 101,075 cases&Prop.~\ref{prop:17-rankone}\\
Rank one, \((3,0,5)\)&Recentering screen; binary projection filter; searches&118 classes; 7,065 searched \(N\)-orbits&Prop.~\ref{prop:17-rankone}\\\bottomrule
\end{tabularx}
\caption{The cases of the proof of Theorem~\ref{thm:main}. The four
rank-two types use the 56 dimension-six shortened-code classes; the
rank-one kernel cases use the dimension-five classes of hull dimension
one, three, and five.}
\label{tab:map}
\end{table}

\subsection{Previous work}

The symplectic representation of stabilizer codes and the MacWilliams and
shadow identities are standard \cite{crss,rains}. Gundersen et al.
\cite{puncturing} describe puncturing directly in symplectic and stabilizer
matrix notation. Complementary symplectic projections are treated in
\cite{anticodes}; related subspace constructions appear in the use of
Goursat's lemma for subsystem codes \cite{goursat} and in
entanglement-assisted error correction \cite{catalytic}.
Our extension criterion uses this linear algebra while retaining the
minimum weights in the original coordinate blocks. Its role here is to
derive the kernel and hull restrictions and to specify the finite
problems used in the nonexistence proof.

Classification by extension and the use of automorphism groups are also
established methods \cite{kaski}. We give completeness checks for the
classifications needed here. Earlier quantum nonexistence proofs
combining structural reductions and finite searches include the
\(\qc{13,5,4}\) result of Bierbrauer et al. \cite{bierbrauer}.
Related geometric methods for purity and local structure are developed
in \cite{geometry}. Singh \cite{singh} announced the nonexistence of
\(\qc{14,3,5}\) codes, including degenerate codes, and Lai \cite{lai},
Section~VII-C1, also proves it.

The semidefinite-programming bounds with rational certificates in
\cite{sdp}, Table~1, leave code-space dimension \(K=8\) possible at
\((n,d)=(17,6)\): their displayed interval is \(4\leq K\leq8\).
Lai \cite{lai} develops refined weight-enumerator and structural tests.
His Table~II records the intermediate candidate counts, and
Section~VII-D leaves seven candidate distributions for
\(\qc{17,3,6}\), so the case remained open. Theorem~\ref{thm:main}
excludes these codes through complete shortened-code classifications and
extension tests.

\subsection{Organization}

Section~\ref{sec:extension} proves the extension criterion, and
Section~\ref{sec:restrictions} derives the kernel and hull identities.
Section~\ref{sec:forced} proves purity, forces the weight-eight word, and
lists the five local types. Section~\ref{sec:classification} proves that
the shortened-code and local-code classifications are complete.
Section~\ref{sec:search} describes the search for extension maps, and
Sections~\ref{sec:rank-two} and \ref{sec:rank-one} exclude the rank-two
and rank-one types. Section~\ref{sec:verification} describes the
computational data and its checks.
